\section{Introduction}
This paper explores how Persistence Mechanics might connect a finite-capacity substrate to cosmological observables. The starting point is the substrate geometry and operational bookkeeping developed in Papers I--V. The cosmological application is less complete: a local capacity count does not by itself determine an expansion history, perturbation spectrum or gravitational field.

We study three proposed applications. First, channel and boundary counts suggest dimensionless primordial and late-time scaling formulas. Second, a reserved coordinate slot supplies an additional contribution to the proposed cosmological load. Third, finite storage and independent-cell counting constrain possible saturation and entropy models. The companion Lean project proves the mathematical consequences of these constructions while leaving their physical identifications explicit.

The numerical comparisons retain the draft's Planck 2018, SH0ES 2022 and DES Year 3 reference inputs. They are fixed benchmarks, not an updated observational review or a joint likelihood analysis. In particular, the early and late $H_0$ labels denote inferences of a present-day rate under stated models; they are not two instantaneous expansion rates at different epochs.

The aim is a dual description of cosmological dynamics from substrate rules. At present the load response, primordial spectrum, baryon-number-changing process, gravitational metric and radiation mechanism require further derivation. Updating the upstream particle calculations also changes some numerical comparisons: the complete CKM ansatz and the older geometric CP formula give different baryogenesis estimates, while a common factor-six response does not repair the proposed neutrino spectrum. Those discrepancies are retained as constraints on the model.
